\section{Mode A1 -- Arrêt en conditions initiales}
\label{secModeA1}

\subsection{Conditions d'entrée}

Le mode A1 est atteint depuis le mode A6, une fois que le robot a rejoint
sa position de référence au-dessus du guichet. Il est également atteint
depuis les modes F1/A2 ou F4, lorsque l'opérateur ramène l'installation à
l'arrêt après un fonctionnement en production ou en mode manuel.

\subsection{Comportement attendu}

\begin{itemize}
    \item \textbf{Robot} : le robot reste immobile, dans sa position de
    référence au-dessus du guichet, tout en demeurant sous puissance, prêt
    à reprendre un mouvement dès qu'un mode de production ou un mode
    manuel sera demandé.
    \item \textbf{Convoyeur / guichet} : le flux des palettes au niveau du
    guichet est libre (passant), le poste n'étant, dans ce mode, le siège
    d'aucune opération de saisie ou de dépose par le robot.
    \item \textbf{IHM} : l'installation est prête pour un fonctionnement
    automatique ou manuel ; cette disponibilité doit être signalée
    clairement à l'opérateur.
\end{itemize}

\subsection{Conditions de sortie}

\begin{itemize}
    \item Vers le mode \textbf{F1} (\emph{Production normale}), sur
    demande de marche automatique de l'opérateur.
    \item Vers le mode \textbf{F4} (\emph{Mode manuel}), sur demande de
    marche manuelle de l'opérateur (ces deux demandes étant mutuellement
    exclusives).
\end{itemize}
